\section{AMC 12 12 2001}
        \begin{enumerate}
        	\item (\textit{AMC 12 12 2001}) The sum of two numbers is $S$. Suppose $3$ is added to each number and then
each of the resulting numbers is doubled. What is the sum of the final two
numbers?


 $\text{(A)}\ 2S + 3\qquad \text{(B)}\ 3S + 2\qquad \text{(C)}\ 3S + 6 \qquad\text{(D)} 2S + 6 \qquad \text{(E)}\ 2S + 12$

 \textbf{Answer: }E



\textbf{Solution}

Suppose the two numbers are $a$ and $b$, with $a+b=S$.Then the desired sum is$2(a+3)+2(b+3)=2(a+b)+12=2S +12$, which is answer $\boxed{\textbf{(E)}}$.

 


	\item (\textit{AMC 12 12 2001}) Let $P(n)$ and $S(n)$ denote the product and the sum, respectively, of the digits
of the integer $n$. For example, $P(23) = 6$ and $S(23) = 5$. Suppose $N$ is a
two-digit number such that $N = P(N)+S(N)$. What is the units digit of $N$?


 $\text{(A)}\ 2\qquad \text{(B)}\ 3\qquad \text{(C)}\ 6\qquad \text{(D)}\ 8\qquad \text{(E)}\ 9$

 \textbf{Answer: }E



\textbf{Solution 1}

Denote $a$ and $b$ as the tens and units digit of $N$, respectively. Then $N = 10a+b$. It follows that $10a+b=ab+a+b$, which implies that $9a=ab$. Since $a\neq0$, $b=9$. So the units digit of $N$ is $\boxed{\textbf{(E) }9}$.

 


	\item (\textit{AMC 12 12 2001}) The state income tax where Kristin lives is levied at the rate of $p\%$ of the first
$\textdollar 28000$ of annual income plus $(p + 2)\%$ of any amount above $\textdollar 28000$. Kristin
noticed that the state income tax she paid amounted to $(p + 0.25)\%$ of her
annual income. What was her annual income?


 $\text{(A)}\ \$28000\qquad \text{(B)}\ \$32000\qquad \text{(C)}\ \$35000\qquad \text{(D)}\ \$42000\qquad \text{(E)}\ \$56000$

 \textbf{Answer: }B



\textbf{Solution 1}

Let $A$, $T$ be Kristin's annual income and the income tax total, respectively. Notice that
 \begin{align*} T &= p\%\cdot28000 + (p + 2)\%\cdot(A - 28000) \\ &= [p\%\cdot28000 + p\%\cdot(A - 28000)] + 2\%\cdot(A - 28000) \\ &= p\%\cdot A + 2\%\cdot(A - 28000) \end{align*}
We are also given that 
 \[T = (p + 0.25)\%\cdot A = p\%\cdot A + 0.25\%\cdot A\]Thus, 
 \[p\%\cdot A + 2\%\cdot(A - 28000) = p\%\cdot A + 0.25\%\cdot A\] 
 \[2\%\cdot(A - 28000) = 0.25\%\cdot A\]Solve for $A$ to obtain $A = \boxed{\textbf{(B) }\textdollar 32000}$.

 Note: You may also assume p = 0 for easy calculation.

 


	\item (\textit{AMC 12 12 2001}) The mean of three numbers is $10$ more than the least of the numbers and $15$
less than the greatest. The median of the three numbers is $5$. What is their
sum?


 $\text{(A)}\ 5\qquad \text{(B)}\ 20\qquad \text{(C)}\ 25\qquad \text{(D)}\ 30\qquad \text{(E)}\ 36$

 \textbf{Answer: }D



\textbf{Solution 1}

Let $m$ be the mean of the three numbers. Then the least of the numbers is $m-10$and the greatest is $m + 15$. The middle of the three numbers is the median, $5$. So$\dfrac{1}{3}[(m-10) + 5 + (m + 15)] = m$, which can be solved to get $m=10$.Hence, the sum of the three numbers is $3\cdot 10 = \boxed{\textbf{(D) }30}$.

 



\textbf{Solution 2}

Say the three numbers are $x$, $y$ and $z$. When we arrange them in ascending order then we can assume $y$ is in the middle therefore $y = 5$. 

 We can also assume that the smallest number is $x$ and the largest number of the three is $z$.Therefore,

 
 \[\frac{x+y+z}{3} = x + 10 = z - 15\]
 \[\frac{x+5+z}{3} = x + 10 = z - 15\]

 Taking up the first equation $\frac{x+5+z}{3} = x + 10$ and simplifying we obtain $z - 2x - 25 = 0$ doing so for the equation $\frac{x+5+z}{3} = z - 15$ we obtain the equation $x - 2z + 50 = 0$

 \[x - 2z + 50 = 2x - 4z + 100\] 

 when solve the above obtained equation and $z - 2x - 25 = 0$ we obtain the values of $z = 25$ and $x = 0$

 Therefore the sum of the three numbers is $25 + 5 + 0 = \boxed{\textbf{(D) }30}$



	\item (\textit{AMC 12 12 2001}) What is the product of all positive odd integers less than 10000?


 $\text{(A)}\ \dfrac{10000!}{(5000!)^2}\qquad \text{(B)}\ \dfrac{10000!}{2^{5000}}\qquad \text{(C)}\ \dfrac{9999!}{2^{5000}}\qquad \text{(D)}\ \dfrac{10000!}{2^{5000} \cdot 5000!}\qquad \text{(E)}\ \dfrac{5000!}{2^{5000}}$

 \textbf{Answer: }D



\textbf{Solution}


 \[1(3)(5) \cdots (9,999) = 9,999!!\]

 We can express the double factorial of an odd number $(2n-1)!!$ in terms of standard factorials using the formula:
 \[(2n-1)!! = \frac{(2n)!}{(2n)!!} = \frac{(2n)!}{2^n(n!)}\]

 Setting $2n = 10,000$, we have $n = 5,000$. Substituting these values into the formula gives:
 \[9,999!! = \frac{10,000!}{2^{5,000}(5,000!)}\]




\textbf{Solution 2 (making the problem easier)}

If you did not see the pattern, then we may solve a easier problem.

 What is the product of all positive odd integers less than $10$?

 1(3)(5)(7)(9) = 945.

 Originally, we had$\textbf{(A)}\ \dfrac{10000!}{(5000!)^2}\qquad \text{(B)}\ \dfrac{10000!}{2^{5000}}\qquad \textbf{(C)}\ \dfrac{9999!}{2^{5000}}\qquad \text{(D)}\ \dfrac{10000!}{2^{5000} \cdot 5000!}\qquad \textbf{(E)}\ \dfrac{5000!}{2^{5000}}$

 but now we have$\textbf{(A)}\ \dfrac{10!}{(5!)^2}\qquad \text{(B)}\ \dfrac{10!}{2^{5}}\qquad \textbf{(C)}\ \dfrac{9!}{2^{5}}\qquad \text{(D)}\ \dfrac{10!}{2^{5} \cdot 5!}\qquad \textbf{(E)}\ \dfrac{5!}{2^{5}}$

 which expression equals 945 

 $\textbf{(A)}\ \dfrac{10!}{(5!)^2}$ = 252 way too small 

 $\textbf{(B)}\ \dfrac{10!}{2^{5}}$ = is way too big, 113400

 $\textbf{(C)}\ \dfrac{9!}{2^{5}}$ = is just 113400 divided by 10(11340), so still too big 

 $\textbf{(D)}\ \dfrac{10!}{2^{5} \cdot 5!}$ = 113400/120= 945, just perfect

 $\textbf{(E)}\ \dfrac{5!}{2^{5}}$ = 3.75 or just too small

 So D is equal to 945, thus the answer is $\boxed{\textbf{(D)} \dfrac{10000!}{2^{5000} \cdot 5000!}}$.

 


	\item (\textit{AMC 12 12 2001}) A telephone number has the form $\text{ABC-DEF-GHIJ}$, where each letter represents
a different digit. The digits in each part of the number are in decreasing
order; that is, $A > B > C$, $D > E > F$, and $G > H > I > J$. Furthermore,
$D$, $E$, and $F$ are consecutive even digits; $G$, $H$, $I$, and $J$ are consecutive odd
digits; and $A + B + C = 9$. Find $A$.


 $\text{(A)}\ 4\qquad \text{(B)}\ 5\qquad \text{(C)}\ 6\qquad \text{(D)}\ 7\qquad \text{(E)}\ 8$

 \textbf{Answer: }E



\textbf{Solution}

We start by noting that there are $10$ letters, meaning there are $10$ digits in total. Listing them all out, we have $0, 1, 2, 3, 4, 5, 6, 7, 8, 9$. Clearly, the most restrictive condition is the consecutive odd digits, so we create casework based on that.

 Case 1: $G$, $H$, $I$, and $J$ are $7$, $5$, $3$, and $1$ respectively. 

 A cursory glance allows us to deduce that the smallest possible sum of $A + B + C$ is $11$ when $D$, $E$, and $F$ are $8$, $6$, and $4$ respectively, so this is out of the question.

 Case 2: $G$, $H$, $I$, and $J$ are $9$, $7$, $5$, and $3$ respectively. 

 A cursory glance allows us to deduce the answer. Clearly, when $D$, $E$, and $F$ are $6$, $4$, and $2$ respectively, $A + B + C$ is $9$ when $A$, $B$, and $C$ are $8$, $1$, and $0$ respectively, giving us a final answer of $\boxed{\textbf{(E)}\ 8}$




\textbf{Solution 2}

The ten letters must all correspond to ten distinct digits, so every digit is used in the telephone number.

 We note that $1$, $3$, $5$, $7$, $9$ are the odd numbers, and possible sequences for $GHIJ$ are either $1$, $3$, $5$, $7$ or $3$, $5$, $7$, $9$.

 $3$, $5$, and $7$ are included in both possible sequences so that we can rule out the possibilities of $3$, $5$, and $7$ for $A$, so $A$ can only be $4$, $6$, or $8$ (Note that we can also rule out $9$ as a possibility since there is no set of $2$ unique digits that could sum with $9$ to $9$). Since every possible sequence of $DEF$ also contains $4$, we can rule out $4$ as well.

 Testing $A$=$6$ means $D$, $E$, $F$ is $0$, $2$, $4$, respectively. However, $6$ and $8$ must then be part of $ABC$, and they already sum to more than $9$. So then we can try $A$=$8$. We have two cases: Where $DEF$ includes $0$, $2$, and $4$, or when $DEF$ includes $2$, $4$, and $6$. If we say the former, $ABC$ includes $6$ as well as $8$, which, again, has a sum over $9$. Therefore, $DEF$ must include $2$, $4$, and $6$, and leaves us with $0$, $1$, and $8$ for $ABC$. This sums to 9, and since the largest digit is $8$, $\boxed{\textbf{(E)}\ 8}$ is our answer.

 


	\item (\textit{AMC 12 12 2001}) A charity sells $140$ benefit tickets for a total of $\textdollar 2001$. Some tickets sell for full price (a whole dollar amount), and the rest sells for half price. How much money is raised by the full-price tickets?


 $\text{(A) }\$782\qquad \text{(B) }\$986\qquad \text{(C) }\$1158\qquad \text{(D) }\$1219\qquad \text{(E) }\$1449$

 \textbf{Answer: }A



\textbf{Solution 1}

Let the cost of the full price ticket be $x$, the number of full-price tickets be $A$, and the number of half-price tickets be $B$

 Let's multiply both sides of the equation that naturally follows by 2. We have

 
 \[2Ax+Bx=4002\]

 And we have $A+B=140\implies B=140-A$

 Plugging in, we get $\implies 2Ax+(140-A)(x)=4002$

 Simplifying, we get $Ax+140x=4002$

 Factoring out the $x$, we get $x(A+140)=4002\implies x=\frac{4002}{A+140}$

 We see that the fraction has to simplify to an integer (the full price is a whole dollar amount)

 Thus, $A+140$ must be a factor of 4002.

 Consider the prime factorization of $4002$: $2\times3\times23\times29$
$A$ must be a positive integer. So, we seek a factor of $4002$ to set equal to $A+140$ so that we get an integer solution for $A$ that is less than $140$. By guess-and-check OR inspection, the appropriate factor is $174$ ($2\times3\times29$), meaning that $A$ has a value of $34$. Plug this into the above equation for $x$ to get $x = 23$.

 Therefore, the price of full tickets out of $2001$ is $23\times34=\boxed{\textbf{(A) }782}$.

 --Edits by Joseph2718 (Reason: Ease of understanding)

 



\textbf{Solution 2}

Let $f$ equal the number of full-price tickets, and let $h$ equal the number of half-price tickets. Additionally, suppose that the price of $f$ is $p$. We are trying to solve for $f \cdot p$.

 Since the total number of tickets sold is $140$, we know that 
 \[f+h=140.\] The sales from full-price tickets ($f \cdot p$) plus the sales from half-price tickets $\Big(\frac{h \cdot p}{2}$, because each hall-price ticket costs $\frac{p}{2}$ dollars$\Big)$ equals $2001.$ Then we can write 
 \[fx + \frac{hx}{2}=2001.\]

 Substituting $h=140-f$ into the second equation, we get 
 \[f \cdot p +\frac{(140-f)p}{2}=f \cdot p+\frac{140p-f\cdot p}{2}=\frac{f\cdot p+140p}{2}=2001.\]

 Multiplying by $2$ and subtracting $140p$ gives us 
 \[f\cdot p=4002-140p.\] 

 Since the problem states that $p$ is a whole number, $140p$ will be some integer multiple of $140$ that ends in a $0$. Thus, $4002-140p$ will end in a $2$. Looking at the answer choices, only $\boxed{\textbf{(A) }782}$ satisfies that condition.

 == Solution 3 This is just a Diophantine Equation!

 Let full-price sales in \# ticket be $f$.  So half-price is just $140-f$.

 Let full-price price be $p$.

 $(140-f)\frac p2+fp=2001$
$(140-f)p+2fp=4002$
$140p+fp=4002$
$(140+f)p=4002$

 Factor $4002=2\times3\times23\times29$.  Note, we need a factor greater than $140$.  It is $2\times3\times29=174$, which makes $140+f=174$, $f=34$, and $p=\frac{4002}{174}=23$.

 $fp=34\times23=782=A$.

 


	\item (\textit{AMC 12 12 2001}) Which of the cones listed below can be formed from a $252^\circ$ sector of a circle of radius $10$ by aligning the two straight sides? 


 
                    \begin{figure}[H]
                        \centering
                        \includegraphics[width=0.5\textwidth]{Problem/AMC_12_12_2001_8_Problem_Figure_1.png}
                    \end{figure}
                    
$\text{(A) A cone with slant height of } 10 \text{ and radius } 6$
$\text{(B) A cone with height of } 10 \text{ and radius } 6$
$\text{(C) A cone with slant height of } 10 \text{ and radius } 7$
$\text{(D) A cone with height of } 10 \text{ and radius } 7$
$\text{(E) A cone with slant height of } 10 \text{ and radius } 8$

 \textbf{Answer: }C



\textbf{Solution 1}


                    \begin{figure}[H]
                        \centering
                        \includegraphics[width=0.5\textwidth]{Solution/AMC_12_12_2001_8_Solution_Figure_1.png}
                    \end{figure}
                    

 The blue lines will be joined together to form a single blue line on the surface of the cone, so $10$ will be the slant height of the cone.

 The red line will form the circumference of the base. We can compute its length and use it to determine the radius.

 The length of the red line is $\dfrac{252^{\circ}}{360^{\circ}}\cdot 2\pi \cdot 10 = \frac{7}{10} \cdot 2\pi \cdot 10 = 14\pi$. This is the circumference of a circle with radius $\frac{14\pi}{2\pi} = 7$.

 Therefore the correct answer is $\boxed{\textbf{(C)} \text{ A cone with slant height of } 10 \text{ and radius } 7}$.

 



\textbf{Solution 2}

As in Solution 1, we deduce that the slant height of the cone is $10$.

 We now recall that the lateral surface area of a cone with base radius $r$ and slant height $l$ is $\pi rl$, while the area of the sector is $\frac{252^{\circ}}{360^{\circ}} \cdot \pi \cdot 10^2 = \frac{7}{10} \cdot \pi \cdot 10^2 = 70\pi$. We must therefore have 
 \[\pi \cdot r \cdot 10 = 70\pi \iff 10r = 70 \iff r = 7,\] so the answer is $\boxed{\textbf{(C)} \text{ A cone with slant height of } 10 \text{ and radius } 7}$.

 -mathlover1205 -formatted and improved by Sevenoptimus

 


	\item (\textit{AMC 12 12 2001}) Let $f$ be a function satisfying $f(xy) = \frac{f(x)}y$ for all positive real numbers $x$ and $y$. If $f(500) =3$, what is the value of $f(600)$?


 $\text{(A)}\ 1 \qquad \text{(B)}\ 2 \qquad \text{(C)}\ \frac52 \qquad \text{(D)}\ 3 \qquad \text{(E)}\ \frac{18}5$

 \textbf{Answer: }C



\textbf{Solution 1}

Letting $x = 500$ and $y = \dfrac{6}{5}$ in the given equation, we get $f(600) = f(500\cdot\frac{6}{5}) = \frac{3}{\left(\frac{6}{5}\right)} = \boxed{\textbf{(C) } \frac{5}{2}}$.

 



\textbf{Solution 2a}

Using the given equation, we obtain 
 \[f(xy) = f(yx) \iff \frac{f(x)}{y} = \frac{f(y)}{x} \iff xf(x) = yf(y),\] for all positive $x$ and $y$. This means $xf(x)$ is constant, i.e. $xf(x) = k$ for some constant $k$, and so the function $f$ must be of the form $f(x) = \frac{k}{x}$. We can now compute $k$ using the given value of $f$: 
 \[f(500) = 3 \iff \frac{k}{500} = 3 \iff k = 1500,\] which yields 
 \[f(600) = \frac{1500}{600} = \boxed{\textbf{(C) } \frac{5}{2}}.\]




\textbf{Solution 2b}

Having determined that $f(x) = \frac{k}{x}$ for some constant $k$, as in Solution 2a, an alternative way to finish the problem is to directly calculate: 

 


	\item (\textit{AMC 12 12 2001}) The plane is tiled by congruent squares and congruent pentagons as indicated. The percent of the plane that is enclosed by the pentagons is closest to


 $\text{(A) }50 \qquad \text{(B) }52  \qquad \text{(C) }54 \qquad \text{(D) }56 \qquad \text{(E) }58$

                    \begin{figure}[H]
                        \centering
                        \includegraphics[width=0.5\textwidth]{Problem/AMC_12_12_2001_10_Problem_Figure_1.png}
                    \end{figure}
                    

 \textbf{Answer: }D



\textbf{Solution}

Consider any single tile:

 
                    \begin{figure}[H]
                        \centering
                        \includegraphics[width=0.5\textwidth]{Solution/AMC_12_12_2001_10_Solution_Figure_1.png}
                    \end{figure}
                    

 If the side of the small square is $a$, then the area of the tile is $\left(3a\right)^2 = 9a^2$, with $4a^2$ covered by squares, and therefore $9a^2-4a^2 = 5a^2$ by pentagons.Hence exactly $\frac{5}{9}$ of any tile is covered by pentagons, and therefore pentagons cover $\frac{5}{9}$ of the plane. When expressed as a percentage, this is $55.\overline{5}\%$, and the closest integer to this value is $\boxed{\textbf{(D) }56}$.

 



\textbf{Solution 2}

There are $6$ smaller square tiles on each side of the bigger square. We are double counting the four corner tiles, however, so we subtract by $4$. Additionally, there are $4$ inner square inside each hole of the two 8 looking figures inside the square. Thus, there are $4$*$2$*$2$ = $16$ inner squares. Thus, there are a total of $36$ squares. Because we are dealing with percentages, and therefore fractions and ratios, we can set the side length to anything we want as long as we maintain consistency throughout the entire figure. For simplicity sake, we set the side length of each small square to $1$

 Thus, the side length of the big square is $9$. Because the amount of squares is $36$, the total area of all the small squares is $36$. The area of the big square is $9$*$9$ = $81$. Hence, $\frac{36}{81}$ or $\frac{4}{9}$ of the entire figure is covered by squares. Thus, $\frac{5}{9}$ is covered by pentagons. Expressed as a decimal, this is $55.\overline{5}\%$. The closest percent value to this is $\boxed{\textbf{(D) }56}$.

 


	\item (\textit{AMC 12 12 2001}) A box contains exactly five chips, three red and two white. Chips are randomly removed one at a time without replacement until all the red chips are drawn or all the white chips are drawn. What is the probability that the last chip drawn is white?


 $\text{(A) }\frac {3}{10} \qquad \text{(B) }\frac {2}{5} \qquad \text{(C) }\frac {1}{2} \qquad \text{(D) }\frac {3}{5} \qquad \text{(E) }\frac {7}{10}$

 \textbf{Answer: }D



\textbf{Solution 1}

Imagine that we draw all the chips in random order, i.e., we do not stop when the last chip of a color is drawn. To draw out all the white chips first, the last chip left must be red, and all previous chips can be drawn in any order. Since there are $3$ red chips, the probability that the last chip of the five is red (and so also the probability that the last chip drawn is white) is $\boxed{\textbf{(D) } \frac {3}{5}}$.

 



\textbf{Solution 2}

Let's assume we don't stop picking until all of the chips are picked. To satisfy this condition, we have to arrange the letters: $W, W, R, R, R$ such that both $W$'s appear in the first $4$. We find the number of ways to arrange the white chips in the first $4$ and divide that by the total ways to choose all the chips. The probability of this occurring is $\dfrac{\dbinom{4}{2}}{\dbinom{5}{2}} = \boxed{\textbf{(D) }\dfrac35}$




\textbf{Solution 3}

The amount of ways to end with a white chip is by having $RRWW, RWW,$ and $WW$. The amount of arrangements for $RRWW$ with $W$ at the end is $3$, the number of arrangements of $RWW$ with $W$ at the end is $2$, and the number of arrangements with $WW$ is just $1$. This gives us $6$ total ways to end with white. Next, the cases to end with a red are $RWRR$, and $RRR$. $RWRR$ gives us $3$ ways and $RRR$ gives us $1$ way. So the number of ways to end with a red is $4$. Thus, our answer is simply $\frac{6}{4+6}$ = $\boxed{\textbf{(D) }\dfrac35}$




\textbf{Solution 4}

First, we figure out the total number of ways to arrange the $5$ chips. There are $\frac{5!}{3!2!} = 10$ ways of arranging them. Now if the last chip is white, all we need to do is to arrange the 4 chips before the last $W$, which are $WWRR$. There are $\frac{4!}{2!2!} = 6$ ways of arranging the first 4 chips. Therefore, the answer is $\frac{6}{10} = \boxed{\textbf{(D) }\dfrac35}$




\textbf{Solution 5}

Test a few possibilities and see that all of the possibilities with white last have probability $\frac{1}{10}$, and there are $6$ possible ways, so it is $\frac{6}{10} = \boxed{\text{(D) }\frac{3}{5}}$.

 


	\item (\textit{AMC 12 12 2001}) How many positive integers not exceeding $2001$ are multiples of $3$ or $4$ but not $5$?


 $\text{(A) }768 \qquad \text{(B) }801 \qquad \text{(C) }934 \qquad \text{(D) }1067 \qquad \text{(E) }1167$

 \textbf{Answer: }B



\textbf{Solution 1}

Out of the numbers $1$ to $12$ four are divisible by $3$ and three by $4$, counting $12$ twice.Hence $6$ out of these $12$ numbers are multiples of $3$ or $4$. 

 The same is obviously true for the numbers $12k+1$ to $12k+12$ for any positive integer $k$.

 Hence out of the numbers $1$ to $60=5\cdot 12$ there are $5\cdot 6=30$ numbers that are divisible by $3$ or $4$.Out of these $30$, the numbers $10$, $15$, $20$, $30$, $40$, $45$ and $60$ are divisible by $5$.Therefore in the set $\{1,\dots,60\}$ there are precisely $30-6=24$ numbers that satisfy all criteria from the problem statement.

 Again, the same is obviously true for the set $\{60k+1,\dots,60k+60\}$ for any positive integer $k$.

 We have $1980/60 = 33$, hence there are $24\cdot 33 = 792$ good numbers among the numbers $1$ to $1980$. At this point we already know that the only answer that is still possible is $\boxed{\textbf{(B)}}$, as we only have $20$ numbers left.

 By examining the remaining $20$ by hand we can easily find out that exactly $9$ of them match all the criteria, giving us $792+9=\boxed{\textbf{(B) }801}$ good numbers.This is correct.

 



\textbf{Solution 2}

We can solve this problem by finding the cases where the number is divisible by $3$ or $4$, then subtract from the cases where none of those cases divide $5$. To solve the ways the numbers divide $3$ or $4$ we find the cases where a number is divisible by $3$ and $4$ as separate cases. We apply the floor function to every case to get $\left\lfloor \frac{2001}{3} \right\rfloor$, $\left\lfloor \frac{2001}{4} \right\rfloor$, and $\left\lfloor \frac{2001}{12} \right\rfloor$. The first two floor functions were for calculating the number of individual cases for $3$ and $4$. The third case was to find any overlapping numbers. The numbers were $667$, $500$, and $166$, respectively. We add the first two terms and subtract the third to get $1001$. The first case is finished.

 The second case is more or less the same, except we are applying $3$ and $4$ to $5$. We must find the cases where the first case over counts multiples of five. Utilizing the floor function again on the fractions $\left\lfloor \frac{2001}{3\cdot5} \right\rfloor$, $\left\lfloor \frac{2001}{4\cdot5} \right\rfloor$, and $\left\lfloor \frac{2001}{3\cdot4\cdot5} \right\rfloor$ yields the numbers $133$, $100$, and $33$. The first two numbers counted all the numbers that were multiples of either four with five or three with five less than $2001$. The third counted the overlapping cases, which we must subtract from the sum of the first two. We do this to reach $200$. Subtracting this number from the original $1001$ numbers procures $\boxed{\textbf{(B)}\ 801}$.

 



\textbf{Solution 3}

First find the number of such integers between 1 and 2000 (inclusive) and then add one to this result because 2001 is a multiple of $3$.

 There are $\frac45\cdot2000=1600$ numbers that are not multiples of $5$.  $\frac23\cdot\frac34\cdot1600=800$ are not multiples of $3$ or $4$, so $800$ numbers are.  $800+1=\boxed{\textbf{(B)}\ 801}$




\textbf{Solution 4}

Take a good-sized sample of consecutive integers; for example, the first $25$ positive integers. Determine that the numbers $3, 4, 6, 8, 9, 12, 16, 18, 21,$ and $24$ exhibit the properties given in the question. $25$ is a divisor of $2000$, so there are $\frac{10}{25}\cdot2000=800$ numbers satisfying the given conditions between $1$ and $2000$. Since $2001$ is a multiple of $3$, add $1$ to $800$ to get $800+1=\boxed{\textbf{(B)}\ 801}$.

 



\textbf{Solution 5}

By PIE, there are $1001$ numbers that are multiples of $3$ or $4$ and less than or equal to $2001$. $80\%$ of them will not be divisible by $5$, and by far the closest number to $80\%$ of $1001$ is $\boxed{\textbf{(B)}\ 801}$.

 



\textbf{Solution 5}

Similar to some of the above solutions. We can divide $2001$ by $3$ and $4$ to find the number of integers divisible by $3$ and $4$. Hence, we find that there are $667$ numbers less than $2001$ that are divisible by $3$, and $500$ numbers that are divisible by $4$. However, we will need to subtract the number of multiples of $15$ from 667 and that of $20$ from $500$, since they're also divisible by 5 which we don't want. There are $133$ + $100$ = $233$ such numbers. Note that during this process, we've subtracted the multiples of $60$ twice because they're divisible by both $15$ and $20$, so we have to add $33$ back to the tally (there are $33$ multiples of $60$ that does not exceed $2001$). Lastly, we have to subtract multiples of both $3$ AND $4$ since we only want multiples of either $3$ or $4$. This is tantamount to subtracting the number of multiples of $12$. And there are $166$ such numbers. Let's now collect our numbers and compute the total: $667$ + $500$ - $133$ - $100$ + $33$ - $166$ = $\boxed{\textbf{(B)}\ 801}$.

 



\textbf{Solution 6 (Venn Diagram)}


                    \begin{figure}[H]
                        \centering
                        \includegraphics[width=0.5\textwidth]{Solution/AMC_12_12_2001_12_Solution_Figure_1.png}
                    \end{figure}
                    

 We have the following Venn diagram (sorry for the bad asymptote). To calculate this Venn diagram we can just use PIE to find the numbers or just work our way out from the inside. We just need to find the numbers that are divisible by $3$ or $4$ but not by $5$, and looking back at the diagram we just have the answer of 
 \[401+133+267 = \boxed{\textbf{(B)}\ 801}\].

 -jb2015007

 



\textbf{Solution 6}

Similar to @above:Let the function $M_{2001}(n)$ return how many multiples of $n$ are there not exceeding $2001$. Then we have that the desired number is:
 \[M_{2001}(3)+M_{2001}(4)-M_{2001}(3\cdot 4)-M_{2001}(3 \cdot 5) - M_{2001}(4 \cdot 5)+M_{2001}(3 \cdot 4 \cdot 5)\]

 Evaluating each of these we get:
 \[667+500-166-133-100+33 = 1100-299 = 801.\]

 Thus, the answer is $\boxed{\textbf{(B)}\ 801}.$

 -ConfidentKoala4

 


 



\textbf{Solution 7}

The smallest common multiple of 3, 4, and 5 is 30. There are 12 integers between 1~30 that are multiples of 3 and 4 but not 5. 2001/30 = 66 R 21. Between 1~21 there are 9 integers that are multiples of 3 and 4 but not 5. 66 * 12 + 9 = $\boxed{\textbf{(B)}\ 801}.$

 - mathlover1205

 - AH2025 (math and number corrections)

 



\textbf{Solution 8}

Let $A$, $B$, and $C$ be the sets of multiples, not exceeding 2001, of $3$, $4$, and $5$, respectively. We desire $|(A\cup B)\backslash C|$. By PIE,\begin{align*}|(A\cup B)\backslash C| &= |A| + |B| - |A \cap B| - |A \cap C| - |B \cap C| + |A \cap B \cap C| \\ &= \lfloor \frac{2001}{3} \rfloor + \lfloor \frac{2001}{4} \rfloor - \lfloor \frac{2001}{3\cdot 4} \rfloor - \lfloor \frac{2001}{3\cdot 5} \rfloor - \lfloor \frac{2001}{4\cdot 5} \rfloor + \lfloor \frac{2001}{3\cdot 4\cdot 5} \rfloor \\&= \boxed{\textbf{(B)}\ 801}\end{align*}
~J.zz.z

 


	\item (\textit{AMC 12 12 2001}) The parabola with equation $y=ax^2+bx+c$ and vertex $(h,k)$ is reflected about the line $y=k$. This results in the parabola with equation $y=dx^2+ex+f$. Which of the following equals $a+b+c+d+e+f$?


 $\text{(A) }2b \qquad \text{(B) }2c \qquad \text{(C) }2a+2b \qquad \text{(D) }2h \qquad \text{(E) }2k$

 \textbf{Answer: }E



\textbf{Solution}

We write $p(x)$ as $a(x-h)^2+k$ (this is possible for any parabola). Then the reflection of $p(x)$ is $q(x) = -a(x-h)^2+k$. Then we find $p(x) + q(x) = 2k$. Since $p(1) = a+b+c$ and $q(1) = d+e+f$, we have $a+b+c+d+e+f = 2k$, so the answer is $\fbox{E}$.

 


	\item (\textit{AMC 12 12 2001}) Given the nine-sided regular polygon $A_1 A_2 A_3 A_4 A_5 A_6 A_7 A_8 A_9$, how many distinct equilateral triangles in the plane of the polygon have at least two vertices in the set $\{A_1,A_2,\dots,A_9\}$?


 $\text{(A) }30 \qquad \text{(B) }36 \qquad \text{(C) }63 \qquad \text{(D) }66 \qquad \text{(E) }72$

 \textbf{Answer: }D



\textbf{Solution}

Each of the $\binom{9}{2} = 36$ pairs of vertices determines $2$ equilateral triangles — one facing towards the center, and one outwards — for a total of $2 \cdot 36 = 72$ triangles. However, the $3$ triangles $A_1A_4A_7$, $A_2A_5A_8$, and $A_3A_6A_9$ are each counted $3$ times (once for each of the $\binom{3}{2} = 3$ possible pairs of vertices, all of which are vertices of the $9$-gon), whereas they should of course only be counted once, resulting in an overcount of $(3-1) \cdot 3 = 6$. Thus, there are $72-6 = \boxed{\text{(D) }66}$ distinct equilateral triangles.

 


	\item (\textit{AMC 12 12 2001}) An insect lives on the surface of a regular tetrahedron with edges of length 1. It wishes to travel on the surface of the tetrahedron from the midpoint of one edge to the midpoint of the opposite edge. What is the length of the shortest such trip? (Note: Two edges of a tetrahedron are opposite if they have no common endpoint.)


 $\text{(A) }\frac {1}{2} \sqrt {3} \qquad \text{(B) }1 \qquad \text{(C) }\sqrt {2} \qquad \text{(D) }\frac {3}{2} \qquad \text{(E) }2$

 \textbf{Answer: }B



\textbf{Solution}

Given any path on the surface, we can unfold the surface into a plane to get a path of the same length in the plane. Consider the net of a tetrahedron in the picture below. A pair of opposite points is marked by dots. It is obvious that in the plane the shortest path is just a segment that connects these two points. By symmetry (as the tetrahedron is regular, so all of its faces are equilateral triangles), its length is the same as the length of the tetrahedron's edge, i.e. $\boxed{\text{(B) }1}$.

 
                    \begin{figure}[H]
                        \centering
                        \includegraphics[width=0.5\textwidth]{Solution/AMC_12_12_2001_15_Solution_Figure_1.png}
                    \end{figure}
                    



	\item (\textit{AMC 12 12 2001}) A spider has one sock and one shoe for each of its eight legs. In how many different orders can the spider put on its socks and shoes, assuming that, on each leg, the sock must be put on before the shoe?


 $\text{(A) }8! \qquad \text{(B) }2^8 \cdot 8! \qquad \text{(C) }(8!)^2 \qquad \text{(D) }\frac {16!}{2^8} \qquad \text{(E) }16!$

 \textbf{Answer: }D



\textbf{Solution 1}

Suppose the spider tries to put on all $2 \cdot 8 = 16$ items in a random order, so that each of the $16!$ possible permutations is equally probable. This means that for any fixed leg, the probability that it will first put on the sock, and only then the shoe, is clearly $\frac{1}{2}$. Hence the probability that it will put on the shoe and sock in the correct order for all its legs is $\left(\frac{1}{2}\right)^8 = \frac{1}{2^{8}}$. Therefore the number of possible permutations is $\boxed{\text{(D) }\frac {16!}{2^8}}$.

 



\textbf{Solution 2}

Each possible order in which the spider puts on its socks and shoes can be uniquely described as a sequence containing two $1$s, two $2$s, ..., and two $8$s, where the first occurrence of number $x$ means that the spider puts the sock onto leg $x$, and the second occurrence means it puts the shoe onto leg $x$ (as it must put on the sock for each leg before the corresponding shoe). 

 Since there are a total of $2 \cdot 8 = 16$ numbers, if they were all unique, the answer would be $16!$. However, since each of the $8$ numbers actually appear twice, there are $2! = 2$ orders for each number that are actually indistinguishable, and so the overcount is by a total factor of $2^8$. This means the answer is $\boxed{\text{(D) }\frac {16!}{2^8}}$.

 



\textbf{Solution 3 (bash using bounding)}

The spider could put on all $8$ socks first, in $8!$ possible ways, and then all $8$ shoes on next, giving $8!$ more possible ways. Of course, this is not the only possibility, so $8! \cdot 8! = (8!)^2$ is a lower bound for the correct answer: the spider could instead, for example, put on $3$ socks and then $1$, $2$, or $3$ shoes, before proceeding to put on more socks.

 In principle, it could put on all $2 \cdot 8 = 16$ items in a random order, but some of these orders would violate the rule that on each leg, the sock must be put on before the shoe, so $16!$ is an upper bound. In other words, answer choices $\text{(C)}$ and $\text{(E)}$ are lower and upper bounds, so the correct answer must lie in between them. Recalling that answer choices on the AMC appear in increasing order, we can deduce, without having to actually calculate it, that the correct answer must be $\boxed{\text{(D) }\frac {16!}{2^8}}$.

 



\textbf{Solution 4}

Name the sock for e.g. leg $1$ as $a1$ and the corresponding shoe as $b1$, and similarly for all the other legs. Then, as the sock for each leg must be put on before the shoe, we observe that the pairs $(a1,b1)$ and $(b1,a1)$ must both correspond to the order $a1$, then $b1$.

 We can now consider constructing an overall order of length $2 \cdot 8 = 16$, starting with $16$ positions that each need to be filled with $ax$ or $bx$ for some integer $x$ between $1$ and $8$. For each of the $8$ legs, we effectively need to choose $2$ of the possible remaining positions, without regard to the order in which they are chosen (since only $1$ order is actually possible, as explained above).

 In other words, we must choose $2$ unordered positions from $16$, then from the remaining $16-2 = 14$, and so on, so the total number of possible orders is 
 \[\binom{16}{2}\binom{14}{2}\dotsm\binom{2}{2} = \frac{16 \cdot 15 \dotsm 2}{2 \cdot 2 \dotsm 2} = \boxed{\text{(D) }\frac {16!}{2^8}}.\]



	\item (\textit{AMC 12 12 2001}) A point $P$ is selected at random from the interior of the pentagon with vertices $A = (0,2)$, $B = (4,0)$, $C = (2 \pi + 1, 0)$, $D = (2 \pi + 1,4)$, and $E=(0,4)$. What is the probability that $\angle APB$ is obtuse?


 $\text{(A) }\frac {1}{5} \qquad \text{(B) }\frac {1}{4} \qquad \text{(C) }\frac {5}{16} \qquad \text{(D) }\frac {3}{8} \qquad \text{(E) }\frac {1}{2}$

 \textbf{Answer: }C



\textbf{Solution}

The angle $APB$ is obtuse if and only if $P$ lies inside the circle with diameter $AB$. (This follows for example from the fact that the inscribed angle is half of the central angle for the same arc.)

 
                    \begin{figure}[H]
                        \centering
                        \includegraphics[width=0.5\textwidth]{Solution/AMC_12_12_2001_17_Solution_Figure_1.png}
                    \end{figure}
                    

 The area of $AFB$ is $[AFB] = \frac {AF\cdot FB}2 = 4$, and the area of $ABCDE$ is $CD\cdot DE - [AFB] = 4\cdot (2\pi+1) - 4 = 8\pi$.

 From the \href{https://artofproblemsolving.com/wiki/index.php?title=Pythagorean\_theorem}{Pythagorean theorem} the length of $AB$ is $\sqrt{2^2 + 4^2} = 2\sqrt{5}$, thus the radius of the circle is $\sqrt{5}$, and the area of the half-circle that is inside $ABCDE$ is $\frac{ 5\pi }2$. 

 Therefore the probability that $APB$ is obtuse is $\frac{ \frac{ 5\pi }2 }{ 8\pi } = \boxed{\text{(C) } \frac 5{16}}$.

 



\textbf{Solution 1.5}

Remember that an obtuse angle is any angle greater than a right angle (90 degrees). This means that if we can figure out all of the points at which $P$ creates a right angle with $A$ and $B$, any points within these bounds will yield an obtuse angle. Notice that by the \href{https://artofproblemsolving.com/wiki/index.php?title=Pythagorean\_theorem}{Pythagorean theorem}, the length of $AB$ is $2\sqrt{5}$. Also notice that when there is a right angle at $P$, $AB$ becomes the hypotenuse. Therefore, the sum of the squares of lengths $AP$ and $BP$ (or the "legs") must equal $AB^2$, or $20$. This can be written as an equation.

 If Point $P$ is situated at $(x, y)$, then the distance from $P$ to $A$ is as follows: $\sqrt{|x-0|^2 + |y-2|^2}$. Accordingly, the distance from $P$ to $B$ is $\sqrt{|x-4|^2 + |y-0|^2}$. Since absolute value and squaring (because we will be squaring the lengths) both eliminate negatives, we can ignore the absolute value and solve. Setting the sum of the squares of these lengths equal to 20 yields $x^2 - 4x + y^2 - 2y = 0$. Completing the square for both $x$ and $y$, we get $(x-2)^2 + (y-1)^2 = 5$, or the equation of a circle with center $(2,1)$ and radius $\sqrt{5}$. Graphing both the pentagon and the circle, we see that the area enclosed is a semicircle. This is$\frac{2\pi*(\sqrt{5})^2}{2}$ or $\frac{5\pi}{2}$.Calculating the area of the pentagon yields $16\pi$. The probability of point P being within the bounds of the semicircle is then$\frac{5\pi}{16\pi}$ $\rightarrow$ $\frac{5}{16}$ or $\boxed{\text{(C)}}$.

 



\textbf{Solution 2}

(Alcumus Solution)

 Since $\angle APB = 90^{\circ}$ if and only if $P$ lies on the semicircle with center $(2,1)$ and radius $\sqrt{5}$, the angle is obtuse if and only if the point $P$ lies inside this semicircle. The semicircle lies entirely inside the pentagon, since the distance, 3, from $(2,1)$ to $\overline{DE}$ is greater than the radius of the circle. Thus the probability that the angle is obtuse is the ratio of the area of the semicircle to the area of the pentagon.

 
                    \begin{figure}[H]
                        \centering
                        \includegraphics[width=0.5\textwidth]{Solution/AMC_12_12_2001_17_Solution_Figure_2.png}
                    \end{figure}
                    Let $O=(0,0)$, $A=(0,2)$, $B=(4,0)$, $C=(2\pi+1,0)$, $D=(2\pi+1,4)$, and $E=(0,4)$. Then the area of the pentagon is
 \[[ABCDE]=[OCDE]-[OAB] = 4\cdot(2\pi+1)-\frac{1}{2}(2\cdot4) = 8\pi,\]and the area of the semicircle is
 \[\frac{1}{2}\pi(\sqrt{5})^2 = \frac{5}{2}\pi.\]The probability is
 \[\frac{\frac{5}{2}\pi}{8\pi} = \boxed{\text{(C) } \frac 5{16}}.\]



	\item (\textit{AMC 12 12 2001}) A circle centered at $A$ with a radius of 1 and a circle centered at $B$ with a radius of 4 are externally tangent. A third circle is tangent to the first two and to one of their common external tangents as shown. The radius of the third circle is


 
                    \begin{figure}[H]
                        \centering
                        \includegraphics[width=0.5\textwidth]{Problem/AMC_12_12_2001_18_Problem_Figure_1.png}
                    \end{figure}
                    
$\text{(A) }\frac {1}{3} \qquad \text{(B) }\frac {2}{5} \qquad \text{(C) }\frac {5}{12} \qquad \text{(D) }\frac {4}{9} \qquad \text{(E) }\frac {1}{2}$

 \textbf{Answer: }D



\textbf{Solution 1}


                    \begin{figure}[H]
                        \centering
                        \includegraphics[width=0.5\textwidth]{Solution/AMC_12_12_2001_18_Solution_Figure_1.png}
                    \end{figure}
                    

 In the triangle $ABC$ we have $AB = 1+4 = 5$ and $BC=4-1 = 3$, thus by the \href{https://artofproblemsolving.com/wiki/index.php?title=Pythagorean\_theorem}{Pythagorean theorem} we have $AC=4$.

 
Let $r$ be the radius of the small circle, and let $s$ be the perpendicular distance from $S$ to $\overline{AC}$. Moreover, the small circle is tangent to both other circles, hence we have $SA=1+r$ and $SB=4+r$.

 We have $SA = \sqrt{s^2 + (1-r)^2}$ and $SB=\sqrt{(4-s)^2 + (4-r)^2}$. Hence we get the following two equations:

 
 \begin{align*} s^2 + (1-r)^2 & = (1+r)^2 \\ (4-s)^2 + (4-r)^2 & = (4+r)^2 \end{align*}


 Simplifying both, we get

 
 \begin{align*} s^2 & = 4r \\ (4-s)^2 & = 16r \end{align*}


 As in our case both $r$ and $s$ are positive, we can divide the second one by the first one to get $\left( \frac{4-s}s \right)^2 = 4$.

 Now there are two possibilities: either $\frac{4-s}s=-2$, or $\frac{4-s}s=2$.

 In the first case clearly $s<0$, which puts the center on the wrong side of $A$, so this is not the correct case. 

 (Note: This case corresponds to the other circle that is tangent to both given circles and the common tangent line. By coincidence, due to the $4:1$ ratio between radii of $A$ and $B$, this circle turns out to have the same radius as circle $B$, with center directly left of center $B$, and tangent to $B$ directly above center $A$.) 

 The second case solves to $s=\frac 43$. We then have $4r = s^2 = \frac {16}9$, hence $r = \boxed{\frac 49}$.

 
More generally, for two large circles of radius $a$ and $b$, the radius $c$ of the small circle is $c = \frac{ab}{\left(\sqrt{a}+\sqrt{b}\right)^2} = \frac{1}{\left(1/\sqrt{a}+1/\sqrt{b}\right)^2}$. 

 Equivalently, we have that $1/\sqrt{c} = 1/\sqrt{a} + 1/\sqrt{b}$.

 



\textbf{Solution 2}

The horizontal line is the equivalent of a circle of curvature $0$, thus we can apply \href{https://artofproblemsolving.com/wiki/index.php?title=Descartes\%27\_Circle\_Formula}{Descartes' Circle Formula}.

 The four circles have curvatures $0, 1, \frac 14$, and $\frac 1r$.

 We have $2\left(0^2+1^2+\frac {1}{4^2}+\frac{1}{r^2}\right)=\left(0+1+\frac 14+\frac 1r\right)^2$

 Simplifying, we get $\frac{34}{16}+\frac{2}{r^2}=\frac{25}{16}+\frac{5}{2r}+\frac{1}{r^2}$

 \[\frac{1}{r^2}-\frac{5}{2r}+\frac{9}{16}=0\]
 \[\frac{16}{r^2}-\frac{40}{r}+9=0\]
 \[\left(\frac{4}{r}-9\right)\left(\frac{4}{r}-1\right)=0\]

 Obviously $r$ cannot equal $4$, therefore $r = \boxed{\frac 49}$.

 



\textbf{Solution 3 (Basically 1 but less complicated)}

As in solution 1, in triangle $ABC$ we have $AB = 1+4 = 5$ and $BC=4-1 = 3$, thus by the Pythagorean theorem or pythagorean triples in general, we have $AC=4$.Let $r$ be the radius. Let $s$ be the perpendicular intersecting point $S$ and line $BC$. $AC=s$ because $s,$ both perpendicular radii, and $AC$ form a rectangle. We just have to find $AC$ in terms of $r$ and solve for $r$ now. From the Pythagorean theorem and subtracting to get lengths, we get $AC=s=4=\sqrt{(r+1)^2 - (1-r)^2} + \sqrt{(r+4)^2 - (4-r)^2}$, which is simply $4=\sqrt{4r}+\sqrt{16r} \implies \sqrt{r}=\frac{2}{3} \implies r= \boxed{\textbf{(D) } \frac{4}{9}}.$



	\item (\textit{AMC 12 12 2001}) The polynomial $P(x)=x^3+ax^2+bx+c$ has the property that the mean of its zeros, the product of its zeros, and the sum of its coefficients are all equal. If the $y$-intercept of the graph of $y=P(x)$ is 2, what is $b$?


 $\text{(A) }-11 \qquad \text{(B) }-10 \qquad \text{(C) }-9 \qquad \text{(D) }1 \qquad \text{(E) }5$

 \textbf{Answer: }A



\textbf{Solution}

We are given $c=2$. So the product of the roots is $-c = -2$ by \href{https://artofproblemsolving.com/wiki/index.php?title=Vieta\%27s\_formulas}{Vieta's formulas}. These also tell us that $\frac{-a}{3}$ is the average of the zeros, so $\frac{-a}3=-2 \implies a = 6$. We are also given that the sum of the coefficients is $-2$, so $1+6+b+2 = -2 \implies b=-11$. So the answer is $\fbox{A}$.

 


	\item (\textit{AMC 12 12 2001}) Points $A = (3,9)$, $B = (1,1)$, $C = (5,3)$, and $D=(a,b)$ lie in the first quadrant and are the vertices of quadrilateral $ABCD$. The quadrilateral formed by joining the midpoints of $\overline{AB}$, $\overline{BC}$, $\overline{CD}$, and $\overline{DA}$ is a square. What is the sum of the coordinates of point $D$?


 $\text{(A) }7 \qquad \text{(B) }9 \qquad \text{(C) }10 \qquad \text{(D) }12 \qquad \text{(E) }16$

 \textbf{Answer: }C



\textbf{Solution}


                    \begin{figure}[H]
                        \centering
                        \includegraphics[width=0.5\textwidth]{Solution/AMC_12_12_2001_20_Solution_Figure_1.png}
                    \end{figure}
                    

 We already know two vertices of the square: $(A+B)/2 = (2,5)$ and $(B+C)/2 = (3,2)$. 

 There are only two possibilities for the other vertices of the square: either they are $(6,3)$ and $(5,6)$, or they are $(0,1)$ and $(-1,4)$. The second case would give us $D$ outside the first quadrant, hence the first case is the correct one. As $(6,3)$ is the midpoint of $CD$, we can compute $D=(7,3)$, and $7+3=\boxed{10}$.

 


	\item (\textit{AMC 12 12 2001}) Four positive integers $a$, $b$, $c$, and $d$ have a product of $8!$ and satisfy:


 
 \begin{align*} ab + a + b & = 524 \\  bc + b + c & = 146 \\  cd + c + d & = 104 \end{align*}



 What is $a-d$?


 $\text{(A) }4 \qquad \text{(B) }6 \qquad \text{(C) }8 \qquad \text{(D) }10 \qquad \text{(E) }12$

 \textbf{Answer: }D



\textbf{Solution}

Using Simon's Favorite Factoring Trick, we can rewrite the three equations as follows:

 
 \begin{align*} (a+1)(b+1) & = 525 \\  (b+1)(c+1) & = 147 \\  (c+1)(d+1) & = 105 \end{align*}


 Let $(e,f,g,h)=(a+1,b+1,c+1,d+1)$. We get:

 
 \begin{align*} ef & = 3\cdot 5\cdot 5\cdot 7 \\  fg & = 3\cdot 7\cdot 7 \\  gh & = 3\cdot 5\cdot 7 \end{align*}


 Clearly $7^2$ divides $fg$. On the other hand, $7^2$ can not divide $f$, as it then would divide $ef$. Similarly, $7^2$ can not divide $g$. Hence $7$ divides both $f$ and $g$. This leaves us with only two cases: $(f,g)=(7,21)$ and $(f,g)=(21,7)$.

 The first case solves to $(e,f,g,h)=(75,7,21,5)$, which gives us $(a,b,c,d)=(74,6,20,4)$, but then $abcd \not= 8!$. We do not need to multiply, it is enough to note e.g. that the left hand side is not divisible by $7$. (Also, a - d equals $70$ in this case, which is way too large to fit the answer choices.)

 The second case solves to $(e,f,g,h)=(25,21,7,15)$, which gives us a valid quadruple $(a,b,c,d)=(24,20,6,14)$, and we have $a-d=24-14 =\boxed{10}$.

 


	\item (\textit{AMC 12 12 2001}) In rectangle $ABCD$, points $F$ and $G$ lie on $AB$ so that $AF=FG=GB$ and $E$ is the midpoint of $\overline{DC}$. Also, $\overline{AC}$ intersects $\overline{EF}$ at $H$ and $\overline{EG}$ at $J$. The area of the rectangle $ABCD$ is $70$. Find the area of triangle $EHJ$.


 $\text{(A) }\frac {5}{2} \qquad \text{(B) }\frac {35}{12} \qquad \text{(C) }3 \qquad \text{(D) }\frac {7}{2} \qquad \text{(E) }\frac {35}{8}$

 \textbf{Answer: }C



\textbf{Solution 1}


                    \begin{figure}[H]
                        \centering
                        \includegraphics[width=0.5\textwidth]{Solution/AMC_12_12_2001_22_Solution_Figure_1.png}
                    \end{figure}
                    

 Note that the triangles $AFH$ and $CEH$ are similar, as they have the same angles. Hence $\frac {AH}{HC} = \frac{AF}{EC} = \frac 23$.

 
	Also, triangles $AGJ$ and $CEJ$ are similar, hence $\frac {AJ}{JC} = \frac {AG}{EC} = \frac 43$.

 
	We can now compute $[EHJ]$ as $[ACD]-[AHD]-[DEH]-[EJC]$. We have:

 
 

 
            \begin{itemize}
        \item $[ACD]=\frac{[ABCD]}2 = 35$.
            \end{itemize}
            
            \begin{itemize}
        \item $[AHD]$ is $2/5$ of $[ACD]$, as these two triangles have the same base $AD$, and $AH$ is $2/5$ of $AC$, therefore also the height from $H$ onto $AD$ is $2/5$ of the height from $C$. Hence $[AHD]=14$.
            \end{itemize}
            
            \begin{itemize}
        \item $[HED]$ is $3/10$ of $[ACD]$, as the base $ED$ is $1/2$ of the base $CD$, and the height from $H$ is $3/5$ of the height from $A$. Hence $[HED]=\frac {21}2$.
            \end{itemize}
            
            \begin{itemize}
        \item $[JEC]$ is $3/14$ of $[ACD]$ for similar reasons, hence $[JEC]=\frac{15}2$.
            \end{itemize}
            
	Therefore $[EHJ]=[ACD]-[AHD]-[DEH]-[EJC]=35-14-\frac {21}2-\frac{15}2 = \boxed{3}$.

 



\textbf{Solution 2}

As in the previous solution, we note the similar triangles and prove that $H$ is in $2/5$ and $J$ in $4/7$ of $AC$. 

 
	We can then compute that $HJ = AC \cdot \left( \frac 47 - \frac 25 \right) = AC \cdot \frac{6}{35}$. 

 
	As $E$ is the midpoint of $CD$, the height from $E$ onto $AC$ is $1/2$ of the height from $D$ onto $AC$. Therefore we have $[EHJ] = \frac{6}{35} \cdot \frac 12 \cdot [ACD] = \frac 3{35} \cdot 35 = \boxed{3}$.

 



\textbf{Solution 3}

Because we see that there are only lines and there is a rectangle, we can coordbash (place this figure on coordinates). Because this is a general figure, we can assume the sides are $7$ and $10$ (or any other two positive real numbers that multiply to 70). We can find $H$ and $J$ by intersecting lines, and then we calculate the area of $EHJ$ using shoelace formula. This yields $\boxed{3}$.

 



\textbf{Solution 4}

Note that triangle $AFH$ is similar to triangle $CEH$ with ratio $\frac{2}{3}$. Similarly, triangle $AGJ$ is similar to triangle $ECJ$ with ratio $\frac{4}{3}$. Thus, if $AC = a$ then we know that $AH = \frac{2}{5}a$ and $JC = \frac{3}{7}a$ meaning $HJ = \frac{6}{35}a$ and thus the ratio of $HJ$ to $JC$ is $\frac{\frac{6}{35}}{\frac{3}{7}} = \frac{2}{5}$ which equals the ratio of the areas of $HJE$ to $JEC$. If $y = AD, x = DC$, then we know that $JEC = \text{(altitude from J to EC)} \cdot EC = \frac{3}{7}y \cdot \frac{1}{2}x \cdot \frac{1}{2}$ and since $xy = 70$ and we want to find $\frac{2}{5}$ of this, we get our answer is $\frac{2}{5} \cdot \frac{3}{7} \cdot \frac{1}{2} \cdot 70 \cdot \frac{1}{2} = \boxed{3}$. -SuperJJ

 



\textbf{Solution 5}

$[CEF] = \frac{[ABCD]}{4} = \frac{35}{2}$
$\triangle CEH \sim \triangle AFH$, $\frac{HE}{HF} = \frac{CE}{AF} = \frac{3}{2}$, $\frac{HE}{EF} = \frac{3}{5}$
$[CEH] = \frac{HE}{EF} \cdot [CEF] = \frac{3}{5} \cdot \frac{35}{2} = \frac{21}{2}$
$\triangle CEH \sim \triangle AFH$, $\frac{AH}{HC} = \frac{AF}{CE} = \frac{2}{3}$, $\frac{AH}{AC} = \frac{2}{5}$, $\frac{CH}{AC} = \frac{3}{5}$
$\triangle CEJ \sim \triangle AGJ$, $\frac{AJ}{JC} = \frac{AG}{CE} = \frac{4}{3}$, $\frac{AJ}{AC} = \frac{4}{7}$
$\frac{HJ}{AC} = \frac{AJ}{AC} - \frac{AH}{AC} = \frac{4}{7} - \frac{2}{5} = \frac{6}{35}$
$\frac{HJ}{CH} = \frac{HJ}{AC} \cdot \frac{AC}{CH} = \frac{6}{35} \cdot \frac{5}{3} = \frac{2}{7}$
$[EHJ] = \frac{HJ}{CH} \cdot [CEH] = \frac{2}{7} \cdot \frac{21}{2} = \boxed{\textbf{(C) }3}$




\textbf{Solution 6 (Mass Points)}

We need one more pair of ratios to fully define our mass point system.  Let's use $\triangle AFH \sim \triangle CEH\implies EH:HF = 3:2$ and now do mass points on $\triangle AEG$:

 
                    \begin{figure}[H]
                        \centering
                        \includegraphics[width=0.5\textwidth]{Solution/AMC_12_12_2001_22_Solution_Figure_2.png}
                    \end{figure}
                    

 Now it's just a standard "Area Reduction by Ratios"™ problem going from:

 
 \[[ABCD]\xrightarrow[]{\frac{1}{2}}[AEB]\xrightarrow[]{\frac{2}{3}}[AEG]\xrightarrow[]{\frac{3}{7}}[AEJ]\xrightarrow[]{\frac{3}{10}}[HEJ]\]

 or,

 
 \[70 \cdot \frac{1}{2}\cdot \frac{2}{3}\cdot \frac{3}{7}\cdot \frac{3}{10} = \boxed{\textbf{(C) }3}\]



	\item (\textit{AMC 12 12 2001}) A polynomial of degree four with leading coefficient 1 and integer coefficients has two real zeros, both of which are integers. Which of the following can also be a zero of the polynomial?


 $\text{(A) }\frac {1 + i \sqrt {11}}{2} \qquad \text{(B) }\frac {1 + i}{2} \qquad \text{(C) }\frac {1}{2} + i \qquad \text{(D) }1 + \frac {i}{2} \qquad \text{(E) }\frac {1 + i \sqrt {13}}{2}$

 \textbf{Answer: }A



\textbf{Solution}

Let the polynomial be $P$ and let the two integer zeros be $z_1$ and $z_2$. We can then write $P(x)=(x-z_1)(x-z_2)(x^2+ax+b)$ for some integers $a$ and $b$.

 If a complex number $p+qi$ with $q\not=0$ is a root of $P$, it must be the root of $x^2+ax+b$, and the other root of $x^2+ax+b$ must be $p-qi$.

 We can then write $x^2+ax+b = (x-p-qi)(x-p+qi) = (x-p)^2 - (qi)^2 = x^2 - 2px + p^2 + q^2$. 

 We can now examine each of the five given complex numbers, and find the one for which the values $-2p$ and $p^2+q^2$ are integers. This is $\boxed{\frac {1 + i \sqrt {11}}{2}}$, for which we have $-2p = -2\cdot\frac 12 = -1$ and $p^2+q^2 = \left( \frac 12 \right)^2 + \left( \frac {\sqrt{11}}2 \right)^2 = \frac 14 + \frac {11}4 = \frac {12}4 = 3$.

 (As an example, the polynomial $x^4 - 2x^3 + 4x^2 - 3x$ has zeroes $0$, $1$, and $\frac {1 \pm i \sqrt {11}}{2}$.)

 



\textbf{Solution 2}

By Vieta, we know that the product of all four zeros of the polynomial equals the constant at the end of the polynomial.We also know that the two imaginary roots are a conjugate pair (I.E if one is a+bi, the other is a-bi). So the two imaginary roots must multiply to give you an integer.Taking the 5 answers into hand, we find that $\boxed{\frac {1 + i \sqrt {11}}{2}}$ is our only integer giving solution.

 Note: I think this solution is not correct. the products of the roots are integers do not mean the product of the two complex roots are integers.

 Note: I believe this solution is correct. We know that the two real solutions are integers and that the final product is an integer. The product of the real solutions multiplied by the product of the complex solutions equals the final product, so we know that the product of the two complex solutions must be integers. We know the two complex solutions are conjugates, so we can test all the answer choices and find that A is the answer. ~A1597412

 Note: No, the solution is not correct. Here's a counterexample: the product of the two integer roots is 4 and the product of the two complex roots is $1/2$. The product of all four roots is $2$, which is clearly an integer.

 FINAL CLARIFICATION: The note above is accurate. This solution is not fully complete in writing, but is valid nonetheless. If the 2 complex roots multiply to a non-integer, the sum of all 4 roots MAY STILL produce an integer. However, if the 2 complex roots multiply to an integer, all 4 roots MUST produce an integer. So we know that $\boxed{ A) \frac {1 + i \sqrt {11}}{2}}$ MUST be the solution.

 There is a simple fix: The complex roots must be algebraic integers, so the product of the root and its conjugate must be an integer.

 



\textbf{Solution 3}

After dividing the polynomial out by $(x-p)$ and $(x-q)$, where p and q are the real roots of the polynomial, we will obtain a quadratic with two complex roots. We can then use the quadratic formula to solve for these complex roots.

 Let's start by using synthetic division to divide $x^4+ax^3+bx^2+cx+d$ by $(x-p)$. Using this method, the quotient becomes $1x^3+(a-p)x^2+(b-pa+p^2)x+(c-bp+ap^2-p^3)+\frac{d-pc+bp^2-ap^3+p^4}{x-q}$. However, we know that there should be no remainder because $(x-p)$ is a factor of the polynomial, so $\frac{d-pc+bp^2-ap^3+p^4}{x-q}$ must equal 0, so $d=-pc+bp^2-ap^3+p^4$. When we divide the expression on the left by -p, we get $c-bp+ap^2-p^3$, so we can replace it in our original synthetic division equation with $\frac {d}{-k}$.

 We then want to synthetically divide $x^3+(a-p)x^2+(b-pa+p^2)x+\frac {d}{-k}$ by the next factor, $(x-q)$. Using the same method as before, we can simplify the quotient to $x^2+(a-p-q)x+\frac{d}{pq}$. Now for the easy part!

 Use the quadratic formula to determine the form of the complex roots.

 $\frac{p+q-a\pm\sqrt{(a-p-q)^2-\frac{4d}{pq}}}{2}$

 Now this is starting to look a lot like answers A and E. Noticing that the real part in each answer choice is $\frac{1}{2}$, $(k+n-a)=1$ and $(a-k-n)^2=1$, and the imaginary part is positive. Furthermore, by Vieta's Formulas, we know that d must be a multiple of p and q, so $\frac{4d}{pq}$ is a multiple of 4. Rearranging the expression, we get:

 $\frac{1+i\sqrt{\frac{4d}{pq}-1}}{2}$

 The radicand therefore must be one less than a multiple of four, which is only the case in $\frac {1 + i \sqrt {11}}{2}$ or $\boxed{A}$.

 



\textbf{Solution 4 (answer choices)}

According to the \href{https://artofproblemsolving.com/wiki/index.php/Complex\_Conjugate\_Root\_Theorem}{Complex Conjugate Root Theorem}, if $a+bi$ is a root then so is $a-bi$.  Then the quadratic $(x-(a+bi))(x-(a-bi))$ is also a polynomial factor of the quartic in question.  Let's check what that expanded quadratic looks like for each answer choice:

 
 \[ \begin{array}{|c|c|} \hline \textbf{Answer Choice} & \textbf{Quadratic Expression} \\ \hline \text{A} & \left(x - \frac{1 + \sqrt{-11}}{2}\right)\left(x - \frac{1 - \sqrt{-11}}{2}\right) = x^2 - x + 3 \\ \text{B} & \left(x - \frac{1 + \sqrt{-1}}{2}\right)\left(x - \frac{1 - \sqrt{-1}}{2}\right) = x^2 - x + \frac{1}{2} \\ \text{C} & \left(x - \frac{1 + 2\sqrt{-1}}{2}\right)\left(x - \frac{1 - 2\sqrt{-1}}{2}\right) = x^2 - x + \frac{5}{4} \\ \text{D} & \left(x - \frac{2 + \sqrt{-1}}{2}\right)\left(x - \frac{2 - \sqrt{-1}}{2}\right) = x^2 - 2x + \frac{5}{4} \\ \text{E} & \left(x - \frac{1 + \sqrt{-13}}{2}\right)\left(x - \frac{1 - \sqrt{-13}}{2}\right) = x^2 - x + \frac{7}{2} \\ \hline \end{array} \]

 Noting that only choice $\textbf{(A)}$ has all integer coefficients, we can clearly see that for the two other integer roots $p$, and $q$, $(x-p)(x-q)(x^2 - x + 3)$ will always give us a polynomial of degree four with leading coefficient 1 and integer coefficients.  Thus the answer must be $\boxed{\textbf{(A)} = \frac {1 + i \sqrt {11}}{2}}$.

 



\textbf{Solution 5}

Suppose $P(x) = (x-z_1)(x-z_2)(x^2+ax+b)$

 Expand P(x): $P(x) = x^4 + ((z_1+z_2)+a)x^3+(-(z_1+z_2)a+b+z_1z_2)x^2+(-(z_1+z_2)+az_1z_2)x$. 

 Notice that the coefficients of $P(x)$ are integers. Comparing the coefficients, we can know that $a$ and $b$ are integers

 Suppose $p+qi$ and $p-qi$ are the two complex roots of p(x).

 Using $z\overline{z} =$|z|$^ 2$ and $b = (p+qi)(p-qi)$ (using Vieta only on $x^2+ax+b=0$)

 so $p^2+q^2$ is an integer

 select (A)

 


 



\textbf{Solution 6 (Middle School math)}

If we already know two integer roots, that means that the polynomial is $P(x)=(x-r_1)(x-r_2)(x^2+bx+c)$.

 We look at the answer choices and think, hmmm. It looks like the quadratic formula to us. 

 We now try to express answer choice A in terms of the quadratic formula. $\frac{1+i\sqrt{11}}{2}=\frac{-b+\sqrt{b^2-4ac}}{2a}$, so $a=1$ (which we already know) and $b=1$. If $b=-11$, and $\sqrt{b^2-4ac}=i\sqrt{11}$, that means that $(-1)^2-4(1)c=-11$, or $c=3$. 

 Choice A just so happens to be the only answer choice that reverse-engineers a degree-2 polynomial with an integral c term. Therefore, $\boxed{\textbf{(A)} = \frac {1 + i \sqrt {11}}{2}}$.

 


	\item (\textit{AMC 12 12 2001}) In $\triangle ABC$, $\angle ABC=45^\circ$. Point $D$ is on $\overline{BC}$ so that $2\cdot BD=CD$ and $\angle DAB=15^\circ$. Find $\angle ACB$.


 $\text{(A) }54^\circ \qquad \text{(B) }60^\circ \qquad \text{(C) }72^\circ \qquad \text{(D) }75^\circ \qquad \text{(E) }90^\circ$

 \textbf{Answer: }D



\textbf{Solution 2 (what happened to solution 1?)}

Draw a good diagram! Now, let's call $BD=t$, so $DC=2t$. Given the rather nice angles of $\angle ABD = 45^\circ$ and $\angle ADC = 60^\circ$ as you can see, let's do trig. Drop an altitude from $A$ to $BC$; call this point $H$. We realize that there is no specific factor of $t$ we can call this just yet, so let $AH=kt$. Notice that in $\triangle{ABH}$ we get $BH=kt$. Using the 60-degree angle in $\triangle{ADH}$, we obtain $DH=\frac{\sqrt{3}}{3}kt$. The comparable ratio is that $BH-DH=t$. If we involve our $k$, we get:

 $kt(\frac{3}{3}-\frac{\sqrt{3}}{3})=t$. Eliminating $t$ and removing radicals from the denominator, we get $k=\frac{3+\sqrt{3}}{2}$. From there, one can easily obtain $HC=3t-kt=\frac{3-\sqrt{3}}{2}t$. Now we finally have a desired ratio. Since $\tan\angle ACH = 2+\sqrt{3}$ upon calculation, we know that $\angle ACH$ can be simplified. Indeed, if you know that $\tan(75)=2+\sqrt{3}$ or even take a minute or two to work out the sine and cosine using $\sin(x)^2+\cos(x)^2=1$, and perhaps the half- or double-angle formulas, you get $\boxed{75^\circ}$.

 



\textbf{Solution 3}

Without loss of generality, we can assume that $BD = 1$ and $CD = 2$. As above, we are able to find that $\angle ADC = 60^\circ$ and $\angle ADB = 120^\circ$.

 Using Law of Sines on triangle $ADB$, we find that 
 \[\frac{1}{\sin15^\circ} = \frac{AD}{\sin 45^\circ} = \frac{AB}{\sin 120^\circ}.\] Since we know that 
 \[\sin 15^\circ = \frac{\sqrt{6}-\sqrt{2}}{4},\] 
 \[\sin 45^\circ = \frac{\sqrt{2}}{2},\] 
 \[\sin 120^\circ = \frac{\sqrt{3}}{2},\] we can compute $AD$ to equal $1+\sqrt{3}$ and $AB$ to be $\frac{3\sqrt{2}+\sqrt{6}}{2}$.

 Next, we apply Law of Cosines to triangle $ADC$ to see that 
 \[AC^2 = (1+\sqrt{3})^2 + 2^2 - (2)(1+\sqrt{3})(2)(\cos 60^\circ).\] Simplifying the right side, we get $AC^2 = 6$, so $AC = \sqrt{6}$.

 Now, we apply Law of Sines to triangle $ABC$ to see that 
 \[\frac{\sqrt{6}}{\sin 45^\circ} = \frac{\frac{3\sqrt{2}+\sqrt{6}}{2}}{\sin ACB}.\] After rearranging and noting that $\sin 45^\circ = \frac{\sqrt{2}}{2}$, we get 
 \[\sin ACB = \frac{\sqrt{6}+3\sqrt{2}}{4\sqrt{3}}.\]

 Dividing the right side by $\sqrt{3}$, we see that 
 \[\sin ACB = \frac{\sqrt{6}+\sqrt{2}}{4},\] so $\angle ACB$ is either $75^\circ$ or $105^\circ$. Since $105^\circ$ is not a choice, we know $\angle ACB = \boxed{75^\circ}$.

 Note that we can also confirm that $\angle ACB \neq 105^\circ$ by computing $\angle CAB$ with Law of Sines.

 



\textbf{Solution 4(FAST)}

Note that $\angle{ADB} = 120$ and $\angle{ADC} = 60$. Seeing these angles makes us think of 30-60-90 triangles. Let $E$ be the foot of the altitude from $A$ to $BC$. This means $\angle{DAE} = 30$ and $\angle{BAE} = 45$. Let $BD = x$ and $DE = y$. This means $AE = y\sqrt{3}$ and since $AE = BE$ we know that $x + y = y\sqrt{3}$. This means $y = \frac{x(\sqrt{3} + 1)}{2}$. This gives $CE = \frac{4x - x(\sqrt{3} + 1)}{2}$. Note that $\tan{\angle{ACE}} = \frac{AE}{CE} = 2 + \sqrt{3}$. Looking that the answer options we see that $\tan{75^\circ} = 2 + \sqrt{3}$. This means the answer is $D$.~coolmath\_2018 6 7

 



\textbf{Solution 5 (Law of Sines)}

$\angle ADB = 120^\circ$, $\angle ADC = 60^\circ$, $\angle DAB = 15^\circ$, let $\angle ACB = \theta$, $\angle DAC = 120^\circ - \theta$

 By the \href{https://artofproblemsolving.com/wiki/index.php?title=Law\_of\_Sines}{Law of Sines}, we have $\frac{CD}{\sin(120^\circ - \theta)} = \frac{AD}{\sin \theta}$
$\space$ $\frac{BD}{\sin15^\circ} = \frac{AD}{\sin45^\circ}$
$\frac{BD}{CD} \cdot \frac{\sin(120^\circ - \theta)}{\sin15^\circ} = \frac{\sin \theta}{\sin45^\circ}$
$\frac{1}{2} \cdot \frac{\sin(120^\circ - \theta)}{\sin15^\circ} = \frac{\sin \theta}{\sin45^\circ}$

 By the \href{https://artofproblemsolving.com/wiki/index.php/Trigonometric\_identities\#Triple-angle\_identities}{Triple-angle Identities}, $\sin 45^\circ = 3\sin15^\circ -4\sin^3 15^\circ$
$\frac{\sin(120^\circ - \theta)}{2} = \frac{\sin \theta}{3\ -4\sin^2 15^\circ}$
$\sin^2 15^\circ = \frac{1 - \cos 30^\circ}{2} = \frac{1 - \frac{\sqrt{3}}{2}}{2} = \frac{2 - \sqrt{3}}{4}$
$\frac{\sin(120^\circ - \theta)}{2} = \frac{\sin \theta}{3 - 4 \cdot \frac{2 - \sqrt{3}}{4}} = \frac{\sin \theta}{1+\sqrt{3}}$
$\frac{\sin \theta}{\sin(120^\circ - \theta)} = \frac{1+\sqrt{3}}{2}$
$\sin(120^\circ - \theta) = \sin120^\circ \cos\theta - \sin\theta \cos120^\circ = \frac{\sqrt{3}}{2} \cdot \cos\theta + \frac{1}{2} \cdot \sin\theta$，so 

 $\frac{\sin \theta}{\frac{\sqrt{3}}{2} \cdot \cos\theta + \frac{1}{2} \cdot \sin\theta} = \frac{1+\sqrt{3}}{2}$
$\frac{\sqrt{3} \cdot \cos\theta + \sin\theta}{2 \sin \theta} = \frac{2}{1+\sqrt{3}}$
$\frac{\sqrt{3}}{2} \cdot \cot \theta = \frac{2}{1+\sqrt{3}} - \frac{1}{2} = \frac{3 - \sqrt{3}}{2 + 2 \sqrt{3}}$
$\cot \theta = \frac{\cos \theta}{\sin \theta} = \frac{\sqrt{3} - 1}{\sqrt{3} + 1}$

 Suppose $\cos \theta = k(\sqrt{3} - 1)$, and $\sin \theta = k(\sqrt{3} + 1)$
$\sin^2 \theta + \cos^2 \theta = 1$, $k^2(\sqrt{3} + 1)^2 + k^2(\sqrt{3} - 1)^2 = 8k^2 = 1$, $k = \frac{1}{2 \sqrt{2}}$
$\sin \theta = \frac{\sqrt{3} + 1}{2 \sqrt{2}}$, $\cos \theta = \frac{\sqrt{3} - 1}{2 \sqrt{2}}$
$\sin 2 \theta = 2 \cdot \frac{\sqrt{3} + 1}{2 \sqrt{2}} \cdot \frac{\sqrt{3} - 1}{2 \sqrt{2}} = \frac{1}{2}$

 Two possible values of $2 \theta$ are $150^\circ$ and $30^\circ$. However we can rule out $30^\circ$ because $\cos 15^\circ$ is positive, while $\cos \theta$ is negative.

 Therefore $2 \theta = 150^\circ$, $\angle ACB = \boxed{\textbf{(D) } 75^\circ }$




\textbf{Solution 6}

For starters, we have $\angle ABD=120^\circ.$ Dropping perpendiculars $\overline{DX}$ and $\overline{CY}$ from $D$ and $C$ to $\overline{AB}$ gives $\angle ADX=120^\circ-45^\circ=75^\circ,$ since $\angle BDX=45^\circ.$

 Without loss of generality, let $BD=1$ and $CD=2.$ This tells us that $BX=DX=\sqrt{2}/2.$ Using trigonometric identities, we find that

 
 \[\tan \angle ADX=\tan 75^\circ=\dfrac{\sqrt 6+\sqrt 2}{\sqrt 6-\sqrt 2}.\]Thus, $AX/DX=\dfrac{\sqrt 6+\sqrt 2}{\sqrt 6-\sqrt 2},$ which gives $AX=\dfrac{2\sqrt 2+\sqrt 6}{2}.$ Thus, $AB=AX+BX=\dfrac{3\sqrt 2+\sqrt 6}{2}.$

 Now, note that $BCY$ is a $45-45-90$ triangle, so $CY=BY=\dfrac{3\sqrt 2}{2}.$ Thus, we have

 
 \[[ABC]=\dfrac{1}{2}\cdot AB\cdot CY=\dfrac{1}{2}\cdot \dfrac{3\sqrt 2+\sqrt 6}{2}\cdot \dfrac{3\sqrt 2}{2}=\dfrac{9+3\sqrt 3}{4}.\]Additionally, note that $AY=AB-BY=\sqrt{6}/2.$ Applying the Pythagorean Theorem to triangle $AYC$ then tells us that $AC=\sqrt{6}.$ By the trigonometric formula for area,

 
 \[[ABC]=\dfrac{1}{2}\cdot BC\cdot AC\cdot \sin \angle ACB=\dfrac{3\sqrt 6}{2}\sin \angle ACB.\]Setting this equal to our other area and solving gives $\sin \angle ACB=\dfrac{\sqrt 6+\sqrt 2}{4},$ so $\angle ACB=\boxed{75^\circ}.$ ~vaporwave

 


	\item (\textit{AMC 12 12 2001}) Consider sequences of positive real numbers of the form $x, 2000, y, \dots$ in which every term after the first is 1 less than the product of its two immediate neighbors. For how many different values of $x$ does the term $2001$ appear somewhere in the sequence?


 $\text{(A) }1 \qquad \text{(B) }2 \qquad \text{(C) }3 \qquad \text{(D) }4 \qquad \text{(E) more than }4$

 \textbf{Answer: }D



\textbf{Solution}

It never hurts to compute a few terms of the sequence in order to get a feel how it looks like. In our case, the definition is that $\forall$ (for all) $n>1:~ a_n = a_{n-1}a_{n+1} - 1$. This can be rewritten as $a_{n+1} = \frac{a_n +1}{a_{n-1}}$. We have $a_1=x$ and $a_2=2000$, and we compute:

 
 \begin{align*} a_3  & = \frac{a_2+1}{a_1} = \frac{2001}x \\ a_4  & = \frac{a_3+1}{a_2} = \frac{ \dfrac{2001}x + 1 }{ 2000 } = \frac{2001 + x}{2000x} \\ a_5 & = \frac{a_4+1}{a_3} = \frac{ \frac{2001 + x}{2000x} + 1 }{ \frac{2001}x } = \frac{ \frac{2001 + 2001x}{2000x} }{ \frac{2001}x } = \frac{1+x}{2000} \\ a_6 & = \frac{a_5+1}{a_4} = \frac{ \frac{1+x}{2000} + 1 }{ \frac{2001 + x}{2000x} } =  \frac{ \frac{2001+x}{2000} }{ \frac{2001 + x}{2000x} } = x \\ a_7 & = \frac{a_6+1}{a_5} = \frac{ x+1 }{ \frac{1+x}{2000} } = 2000 \end{align*}


 At this point we see that the sequence will become periodic: we have $a_6=a_1$, $a_7=a_2$, and each subsequent term is uniquely determined by the previous two.

 Hence if $2001$ appears, it has to be one of $a_1$ to $a_5$. As $a_2=2000$, we only have four possibilities left. Clearly $a_1=2001$ for $x=2001$, and $a_3=2001$ for $x=1$. The equation $a_4=2001$ solves to $x = \frac{2001}{2000\cdot 2001 - 1}$, and the equation $a_5=2001$ to $x=2000\cdot 2001 - 1$.

 No two values of $x$ we just computed are equal, and therefore there are $\boxed{4}$ different values of $x$ for which the sequence contains the value $2001$.

 \textbf{Remark} In general, notice that the sequence defined by $a_1=a, a_2=b$, and $a_{n+1} = \frac{a_n +1}{a_{n-1}}$ for $n\ge 2$ is periodic.

 


\end{enumerate}